\documentclass[preprint,12pt]{elsarticle}

\usepackage[UTF8]{ctex}
\usepackage{amsmath,amssymb,amsfonts}
\usepackage{booktabs}
\usepackage{graphicx}
\usepackage{hyperref}
\usepackage{algorithm}
\usepackage{algorithmic}

\journal{预印本}

\begin{document}

\begin{frontmatter}

\title{注意力重建正则化用于视觉 Transformer 微调}

\author[inst1]{第一作者\corref{cor1}}
\ead{first.author@example.com}
\author[inst1]{第二作者}

\cortext[cor1]{通讯作者}
\affiliation[inst1]{organization={XXX大学 XXX系}, city={城市}, country={国家}}

\begin{abstract}
基于预训练视觉 Transformer 的迁移学习已成为视觉识别任务的主流范式。然而，标准微调仅优化特定任务的分类目标，使得自注意力权重矩阵中编码的丰富结构信息不受约束。这可能导致在适应过程中，已学习的注意力模式发生过度扭曲，尤其是在容易过拟合的小规模目标数据集上。

本文提出注意力重建正则化（ARR），一种轻量级辅助损失，在微调过程中显式鼓励注意力权重的结构一致性。具体而言，在选定的 Transformer 块中，我们通过 Tikhonov 正则化最小二乘法从注意力输出和值矩阵重建注意力权重矩阵，使用基于 KKT 条件的闭式算法将解投影到概率单纯形上，并惩罚重建注意力分布与原始注意力分布之间的差异。

ARR 不引入任何额外参数，仅在训练时使用。在七个细粒度和纹理识别基准上使用 Swin Transformer-Base 进行的大量实验表明，ARR 在七个数据集中的五个上持续提升准确率，在 DTD 上最高提升 $+1.55\%$，在 FGVC Aircraft 上最高提升 $+1.45\%$，同时计算开销可忽略不计。我们还提供了关于损失权重、正则化放置位置和预热调度的全面消融研究，为在迁移学习中部署基于注意力的正则化提供了实用指南。
\end{abstract}

\begin{keyword}
视觉 Transformer \sep 注意力正则化 \sep 迁移学习 \sep 微调 \sep Swin Transformer
\end{keyword}

\end{frontmatter}

%% ============================================================
\section{引言}
%% ============================================================

视觉 Transformer（ViT）\cite{dosovitskiy2021image} 及其层次化变体如 Swin Transformer \cite{liu2021swin} 已在广泛的视觉识别任务中建立了新的最先进结果。将这些模型部署到下游数据集的主流方法是迁移学习：在大规模源域（如 ImageNet-22K）上预训练的骨干网络，使用标准交叉熵损失在目标任务上进行端到端微调。这种范式已被证明非常有效，但它没有提供显式机制来保留或调节预训练模型所获得的内部结构知识——最值得注意的是自注意力层中的注意力权重分布。

自注意力是 Transformer 的核心计算原语。在每个注意力头中，查询-键内积产生一个权重矩阵 $S \in \mathbb{R}^{N \times N}$（经 softmax 归一化后），控制如何从值 token 中聚合信息。在数百万张图像上的预训练赋予了这些注意力模式丰富的、可泛化的先验：空间局部性、物体-部件关系和纹理敏感的感受野 \cite{raghu2021vision,park2022vision}。然而在微调过程中，注意力权重可以在分类梯度的驱动下以任意方式改变。在小规模目标数据集上——这在细粒度识别、纹理分类和医学影像中很常见——这种无约束的适应可能扭曲已学习的注意力结构，导致过拟合和泛化能力下降。

已有多条研究线路试图正则化微调过程。权重空间正则化方法，如 L2-SP \cite{li2018explicit} 和 DELTA \cite{li2019delta}，惩罚参数偏离预训练初始化的程度。知识蒸馏方法 \cite{hinton2015distilling} 使用预训练模型作为教师来约束微调模型的输出分布。最近，注意力级别的监督已在自蒸馏 \cite{zhang2019your} 和注意力迁移 \cite{zagoruyko2017paying} 的背景下被探索，但这些方法通常需要冻结的教师网络、引入额外参数，或在特征图级别而非直接在注意力权重矩阵上操作。

本文采取不同的视角。我们不是将微调后的注意力权重与冻结的预训练副本进行比较（这会使内存占用翻倍），而是提出通过注意力权重、值矩阵和注意力输出之间���代数关系推导出的\textbf{自一致性约束}来正则化注意力。具体而言，给定选定 Transformer 块的注意力输出 $O = SV$，我们通过 Tikhonov 正则化最小二乘法求解逆问题 $O \approx \tilde{S}V$ 来重建注意力权重矩阵 $\tilde{S}$，并使用闭式 KKT 算法 \cite{duchi2008efficient} 将解投影到概率单纯形上。重建损失
\begin{equation}
\mathcal{L}_{\mathrm{rec}} = \frac{\|\tilde{S} - S_{\mathrm{gt}}\|_F^2}{n \cdot \mathbb{E}[S_{\mathrm{gt}}^2]}
\end{equation}
衡量注意力输出在约束重建权重形成有效概率分布的条件下，能多忠实地被值矩阵"解释"。

我们将此方法称为\textbf{注意力重建正则化（ARR）}。它具有以下特性：
\begin{itemize}
    \item \textbf{零额外参数。} ARR 完全在注意力模块内的现有中间表示（$O$, $V$, $S$）上操作，不需要教师网络、投影头或辅助模块。
    \item \textbf{仅在训练时使用。} 重建损失仅在训练时计算，推理时不增加任何成本。
    \item \textbf{灵活放置。} 正则化可以应用于任意 Transformer 块子集。
    \item \textbf{与现有技术互补。} ARR 正则化注意力结构，与权重衰减、dropout 或输出级蒸馏正交。
\end{itemize}

本文的贡献总结如下：
\begin{enumerate}
    \item 我们提出了一种新颖的注意力级别自一致性正则化方法，基于通过 Tikhonov 正则化和单纯形投影从输出-值线性系统重建注意力权重。
    \item 我们在七个基准上进行了大量实验，并提供了关于损失权重选择（$\alpha = 0.1$）、正则化放置（倒数第二阶段优先）和预热调度（对训练样本少于 4{,}000 的小数据集有益）的实用指南。
    \item 我们坦诚分析了失败模式——准确率饱和和类内变异性——划定了注意力正则化在迁移学习中的适用边界。
\end{enumerate}

%% ============================================================
\section{相关工作}
%% ============================================================

我们将本工作定位在三条研究线路的交汇处：视觉 Transformer 架构、迁移学习正则化策略和注意力级别监督。

\subsection{视觉 Transformer}

Dosovitskiy et al.~\cite{dosovitskiy2021image} 提出了视觉 Transformer（ViT），将标准 Transformer 编码器 \cite{vaswani2017attention} 应用于图像块序列。DeiT \cite{touvron2021training} 证明了精心设计的训练方案可以使 ViT 实现数据高效。Swin Transformer \cite{liu2021swin} 引入了基于移位窗口的自注意力的层次化架构，将全局注意力的二次复杂度降低为图像大小的线性复杂度。进一步的变体包括 PVT \cite{wang2021pyramid}、Twins \cite{chu2021twins} 和 CSWin \cite{dong2022cswin}。我们的工作基于 Swin Transformer，但所提出的正则化是架构无关的：它适用于任何计算 $O = SV$ 的 Transformer 块。

\subsection{迁移学习正则化}

通过微调预训练骨干网络进行迁移学习是有限目标数据下视觉识别的事实标准 \cite{zhuang2020comprehensive}。

\textbf{参数空间正则化。} L2-SP \cite{li2018explicit} 惩罚微调权重与预训练权重之间的 $L_2$ 距离。DELTA \cite{li2019delta} 通过中间激活导出的注意力图在特征空间中扩展了这一思想。BSS \cite{chen2019catastrophic} 抑制分类器权重矩阵的小奇异值以缓解负迁移。

\textbf{参数高效微调。} Adapter \cite{houlsby2019parameter}、LoRA \cite{hu2022lora} 和视觉提示调优 VPT \cite{jia2022visual} 添加不到 $1\%$ 的额外参数，同时达到有竞争力的准确率。

\textbf{输出空间正则化。} 知识蒸馏 \cite{hinton2015distilling} 使用教师网络的软化输出概率来指导学生。R-Drop \cite{liang2021rdrop} 正则化两次不同 dropout 掩码前向传播之间的输出一致性。标签平滑 \cite{szegedy2016rethinking} 和数据增强策略如 CutMix \cite{yun2019cutmix} 和 Mixup \cite{zhang2018mixup} 提供互补的正则化。我们的方法在注意力权重级别进行正则化，与参数漂移、输出 logit 或数据扰��正交。

\subsection{注意力监督与蒸馏}

\textbf{注意力迁移。} Zagoruyko and Komodakis \cite{zagoruyko2017paying} 提出将空间注意力图从教师 CNN 迁移到学生 CNN。Wang et al.~\cite{wang2022minivit} 将此扩展到 Transformer 中的自注意力。这些方法需要冻结的教师网络，使内存占用翻倍。

\textbf{自蒸馏。} Zhang et al.~\cite{zhang2019your} 表明网络可以从自身更深层向更浅层蒸馏知识。DearKD \cite{chen2022dearkd} 为 ViT 提出了两阶段蒸馏方案。SEED \cite{fang2021seed} 引入了自监督蒸馏用于预训练。

\textbf{注意力正则化。} DropKey \cite{li2023dropkey} 在 softmax 之前对注意力 logit 应用 dropout。Re-attention \cite{zhou2021deepvit} 引入可学习变换在头之间混合注意力图。Gong et al.~\cite{gong2021improve} 正则化注意力熵以防止退化分布。相比之下，ARR 完全不改变模型的前向传播；它添加一个纯辅助损失，源自从输出重建注意力权重的逆问题。

\textbf{神经网络中的逆问题。} 我们的公式——求解 $O \approx \tilde{S}V$ 中的 $\tilde{S}$——是一个约束线性逆问题。Tikhonov 正则化 \cite{tikhonov1977solutions} 确保数值稳定性，通过 KKT 条件的单纯形投影遵循 Duchi et al.~\cite{duchi2008efficient} 的高效算法。据我们所知，这是首次将注意力正则化表述为逆重建问题并将其应用于视觉 Transformer 微调的工作。

%% ============================================================
\section{方法}
%% ============================================================

\subsection{预备知识：基于窗口的自注意力}

Swin Transformer \cite{liu2021swin} 将每个特征图划分为大小为 $M \times M$（默认 $M = 7$）的非重叠局部窗口，并在每个窗口内独立计算自注意力。令 $N = M^2$ 为每个窗口的 token 数。对于给定 Transformer 块中的每个注意力头 $h$，输入 token $X \in \mathbb{R}^{N \times C}$ 被投影为查询、键和值：
\begin{equation}
Q = XW_Q, \quad K = XW_K, \quad V = XW_V,
\end{equation}
其中 $Q, K, V \in \mathbb{R}^{N \times d}$，$d = C / n_H$ 是每头维度。注意力权重矩阵和输出计算为：
\begin{equation}
S = \mathrm{softmax}\!\left(\frac{QK^\top}{\sqrt{d}} + B\right), \quad O = SV,
\end{equation}
其中 $B$ 是可学习的相对位置偏置，$S \in \mathbb{R}^{N \times N}$ 满足 $S \geq 0$, $S\mathbf{1} = \mathbf{1}$（每行位于概率单纯形 $\Delta_{N-1}$ 上）。

在 Swin-Base 配置中，四个阶段的深度为 $(2, 2, 18, 2)$，注意力头数为 $(4, 8, 16, 32)$，嵌入维度为 $(128, 256, 512, 1024)$。总参数量为 86.8M。

\subsection{注意力权重重建}

ARR 的核心观察是关系 $O = SV$ 定义了一个关于未知 $S$ 的线性系统，给定观测到的输出 $O$ 和值矩阵 $V$。实际中，$V \in \mathbb{R}^{N \times d}$ 且 $N > d$（例如 $N = 49$, $d = 32$），因此系统对 $S$ 的每行是欠定的且病态的。因此我们采用 Tikhonov 正则化最小二乘：
\begin{equation}
\tilde{S} = \arg\min_{S'} \|O - S'V\|_F^2 + \lambda \|S'\|_F^2,
\end{equation}
其闭式解为：
\begin{equation}\label{eq:tikhonov}
\tilde{S} = OV^\top \left(VV^\top + (\lambda + \varepsilon)I\right)^{-1},
\end{equation}
其中 $\lambda = 0.05$ 是 Tikhonov 正则化系数，$\varepsilon = 10^{-3}$ 是数值稳定项。矩阵求逆在 $N \times N$ 矩阵（窗口大小 7 时为 $49 \times 49$）���操作，计算开销很小。

\textbf{梯度流。} $\tilde{S}$ 从 $O$ 和 $V$ 计算得到，两者都是模型参数的可微函数。真值目标 $S_{\mathrm{gt}} = S.\mathrm{detach}()$ 被视为固定监督信号（梯度不通过 $S_{\mathrm{gt}}$ 流动）。

\subsection{通过 KKT 条件的单纯形投影}

公式~\eqref{eq:tikhonov} 的无约束解 $\tilde{S}$ 通常不满足概率单纯形约束。为强制这些约束，我们将每行 $\tilde{s}_i \in \mathbb{R}^N$ 投影到单纯形 $\Delta_{N-1}$ 上：
\begin{equation}
s_i^* = \arg\min_{s \in \Delta_{N-1}} \|s - \tilde{s}_i\|_2^2.
\end{equation}
KKT 驻点条件给出 $s_j^* = \max(\tilde{s}_{i,j} - \theta, 0)$，其中阈值 $\theta$ 由等式约束 $\sum_j s_j^* = 1$ 确定。

完整算法对每行执行如下步骤：
\begin{enumerate}
    \item 按降序排列条目：$u_1 \geq u_2 \geq \cdots \geq u_N$。
    \item 计算累积和 $C_k = \sum_{j=1}^{k} u_j$。
    \item 找到 $\rho = \max\{k : u_k > (C_k - 1)/k\}$。
    \item 设置 $\theta = (C_\rho - 1)/\rho$。
    \item 返回 $s_j^* = \max(\tilde{s}_{i,j} - \theta, 0)$。
\end{enumerate}
每行复杂度为 $O(N \log N)$。投影是完全可微的。

\subsection{归一化重建损失}

投影后，我们得到 $\tilde{S} \in \Delta_{N-1}$（每行），并与真值注意力权重 $S_{\mathrm{gt}} = S.\mathrm{detach}()$ 比较。我们采用相对 MSE：
\begin{equation}
\mathcal{L}_{\mathrm{rec}} = \frac{\|\tilde{S} - S_{\mathrm{gt}}\|_F^2}{n \cdot \mathbb{E}[S_{\mathrm{gt}}^2]},
\end{equation}
其中 $n$ 是总元素数，$\mathbb{E}[S_{\mathrm{gt}}^2] = \|S_{\mathrm{gt}}\|_F^2 / n$。此归一化确保 $\mathcal{L}_{\mathrm{rec}}$ 是尺度不变的。

当选择多个 Transformer 块进行正则化时，各块损失取平均：
\begin{equation}
\bar{\mathcal{L}}_{\mathrm{rec}} = \frac{1}{|\mathcal{B}|} \sum_{b \in \mathcal{B}} \mathcal{L}_{\mathrm{rec}}^{(b)}.
\end{equation}

\subsection{总损失与预热调度}

总训练损失将标准交叉熵分类损失与重建正则化结合：
\begin{equation}
\mathcal{L}_{\mathrm{total}} = (1 - \alpha)\,\mathcal{L}_{\mathrm{cls}} + \alpha\,\bar{\mathcal{L}}_{\mathrm{rec}},
\end{equation}
其中 $\alpha \in [0, 1)$ 控制正则化强度。实验确定 $\alpha = 0.1$ 是跨数据集的稳健默认值。

\textbf{块选择。} 重建损失可应用于任意 Transformer 块子集。我们考虑五种配置：\texttt{last\_last}（Stage 4 最后一个块）、\texttt{second\_last}（Stage 3 最后一个块）、\texttt{last\_all}（Stage 4 所有块）、\texttt{all\_last}（每个阶段最后一个块）、\texttt{last\_second}（Stage 4 倒数第二个块）。

\textbf{预热调度。} 我们引入余弦预热调度：
\begin{equation}
\alpha(t) = \frac{\alpha_{\max}}{2}\left(1 - \cos\frac{\pi t}{T}\right),
\end{equation}
其中 $t$ 是当前 epoch，$T$ 是总 epoch 数。预热对训练样本少于约 $4{,}000$ 的数据集特别有益。

%% ============================================================
\section{实验}
%% ============================================================

\subsection{实验设置}

\textbf{骨干网络。} 使用在 ImageNet-22K 上预训练的 Swin Transformer-Base \cite{liu2021swin}（86.8M 参数，patch 大小 4，窗口大小 7，输入分辨率 $224 \times 224$）。

\textbf{数据集。} 在七个基准上评估（表~\ref{tab:datasets}）。

\begin{table}[t]
\centering
\caption{数据集统计信息。}\label{tab:datasets}
\begin{tabular}{lcccc}
\toprule
数据集 & 类别数 & 训练集 & 测试集 & 领域 \\
\midrule
Oxford-IIIT Pet & 37 & 3,680 & 3,669 & 宠物品种 \\
DTD & 47 & 3,760 & 1,880 & 纹理 \\
CUB-200-2011 & 200 & 5,994 & 5,794 & 鸟类物种 \\
Stanford Cars & 196 & 8,144 & 8,041 & 汽车型号 \\
FGVC Aircraft & 100 & 6,667 & 3,333 & 飞机变体 \\
Flowers-102 & 102 & 1,020 & 6,149 & 花卉物种 \\
Stanford Dogs & 120 & 12,000 & 8,580 & 犬种 \\
\bottomrule
\end{tabular}
\end{table}

\textbf{训练细节。} 所有模型使用 AdamW 优化器，学习率 $10^{-4}$，权重衰减 0.05。使用余弦退火学习率调度，训练 30 个 epoch。每 GPU 批大小 64，在 $2\times$ NVIDIA A100-40GB GPU 上使用 DDP 训练。

\textbf{默认 ARR 配置。} $\alpha = 0.1$，\texttt{recon\_blocks = second\_last}，Tikhonov $\lambda = 0.05$，$\varepsilon = 10^{-3}$，启用 KKT 单纯形投影，归一化 MSE，无预热调度。

\subsection{主要结果}

\begin{table}[t]
\centering
\caption{主要结果：分类准确率（\%）。粗体表示相对基线的提升。}\label{tab:main}
\begin{tabular}{lccccc}
\toprule
数据集 & 基线 & ARR（默认） & $\Delta$ & 最佳 ARR & 最佳 $\Delta$ \\
\midrule
Oxford-IIIT Pet & 93.24 & 93.84 & \textbf{+0.60} & 94.01 & \textbf{+0.77} \\
DTD & 81.06 & 82.35 & \textbf{+1.29} & 82.61 & \textbf{+1.55} \\
CUB-200-2011 & 89.58 & 89.71 & \textbf{+0.13} & 89.71 & \textbf{+0.13} \\
Stanford Cars & 89.88 & 90.71 & \textbf{+0.83} & 90.71 & \textbf{+0.83} \\
FGVC Aircraft & 73.56 & 74.60 & \textbf{+1.04} & 75.01 & \textbf{+1.45} \\
Flowers-102 & 99.41 & 99.31 & $-0.10$ & 99.41 & $0.00$ \\
Stanford Dogs & 86.55 & 86.12 & $-0.43$ & 86.55 & $0.00$ \\
\bottomrule
\end{tabular}
\end{table}

表~\ref{tab:main} 报告了基线和 ARR 在所有七个数据集上的分类准确率。在默认配置下，ARR 在七个数据集中的五个上提升了准确率。最大增益出现在 DTD（$+1.29\%$）和 FGVC Aircraft（$+1.04\%$）。Stanford Dogs 是唯一有明显负面效果的数据集（$-0.43\%$），我们在第~\ref{sec:discussion} 节分析原因。

\subsection{消融实验：损失权重 $\alpha$}

\begin{table}[t]
\centering
\caption{损失权重 $\alpha$ 对 DTD 的影响（\texttt{last\_last}，seed = 42）。}\label{tab:alpha}
\begin{tabular}{lccccc}
\toprule
$\alpha$ & Ep.~5 & Ep.~15 & Ep.~30 & 最佳 & $\Delta$ \\
\midrule
0（基线） & 77.98 & 79.04 & 81.06 & 81.06 & --- \\
0.001 & 76.28 & 79.63 & 81.76 & 81.97 & +0.91 \\
0.05 & 78.78 & 80.00 & 80.64 & 80.90 & $-0.16$ \\
0.1 & 80.11 & 82.13 & 82.29 & \textbf{82.34} & \textbf{+1.28} \\
0.2 & 79.04 & 79.36 & 81.17 & 81.33 & +0.27 \\
0.5 & 77.07 & 79.26 & 80.53 & 81.17 & +0.11 \\
\bottomrule
\end{tabular}
\end{table}

结果呈现清晰的单峰模式：$\alpha = 0.1$ 最优（表~\ref{tab:alpha}）。

\subsection{消融实验：块放置}

\begin{table}[t]
\centering
\caption{块放置的影响（$\alpha = 0.1$，seed = 42）。}\label{tab:placement}
\begin{tabular}{lcccc}
\toprule
\texttt{recon\_blocks} & DTD (\%) & $\Delta$ & Pet (\%) & $\Delta$ \\
\midrule
基线 & 81.06 & --- & 93.24 & --- \\
\texttt{last\_last} & 82.34 & +1.28 & 93.51 & +0.27 \\
\texttt{second\_last} & 82.35 & +1.29 & 93.84 & +0.60 \\
\texttt{last\_all} & \textbf{82.39} & \textbf{+1.33} & \textbf{94.01} & \textbf{+0.77} \\
\texttt{all\_last} & 82.07 & +1.01 & 93.95 & +0.71 \\
\bottomrule
\end{tabular}
\end{table}

\subsection{消融实验：预热调度}

\begin{table}[t]
\centering
\caption{预热调度的影响（\texttt{second\_last}，$\alpha_{\max} = 0.1$）。}\label{tab:warmup}
\begin{tabular}{lcccc}
\toprule
数据集 & 基线 & 固定 $\alpha$ & 预热 & 预热 vs 固定 \\
\midrule
Oxford-IIIT Pet & 93.24 & 93.84 & 93.90 & +0.06 \\
DTD & 81.06 & 82.35 & \textbf{82.61} & \textbf{+0.26} \\
CUB-200-2011 & 89.58 & 89.71 & 89.32 & $-0.39$ \\
Stanford Cars & 89.88 & 90.71 & 90.49 & $-0.22$ \\
Flowers-102 & 99.41 & 99.31 & 99.31 & $0.00$ \\
Stanford Dogs & 86.55 & 86.12 & 86.29 & +0.17 \\
\bottomrule
\end{tabular}
\end{table}

预热对小数据集有益（表~\ref{tab:warmup}），建议对训练样本少于约 $4{,}000$ 的数据集使用。

\subsection{计算开销}

\begin{table}[t]
\centering
\caption{ARR 的训练计算开销。推理不受影响。}\label{tab:overhead}
\begin{tabular}{lcc}
\toprule
 & 基线 & ARR (\texttt{second\_last}) \\
\midrule
峰值 GPU 内存（每 GPU） & $\sim$11\,GB & $\sim$14.7\,GB (+3.5\,GB) \\
训练吞吐量 & $1.0\times$ & $\sim 0.82\times$ ($-18\%$) \\
额外参数 & 0 & 0 \\
推理延迟 & $1.0\times$ & $1.0\times$ \\
\bottomrule
\end{tabular}
\end{table}

%% ============================================================
\section{讨论}\label{sec:discussion}
%% ============================================================

\subsection{ARR 为什么有效？}

重建损失作为注意力一致性正则化器：它惩罚那些输出无法被忠实地解释为值 token 凸组合的注意力头。高重建误差 $\|\tilde{S} - S_{\mathrm{gt}}\|^2$ 意味着从 $V$ 到 $O$ 的映射变得病态。通过最小化此差异，ARR 隐式约束注意力分布的秩和平滑度，阻止因在小数据集上过拟合而产生的退化模式。

这种机制对纹理和细粒度识别任务特别有效。在 DTD（$+1.55\%$）、FGVC Aircraft（$+1.45\%$）和 Stanford Cars（$+0.83\%$）上，下游任务恰好需要预训练注意力中编码的局部纹理模式和部件级空间关系。

\subsection{ARR 何时失效？}

\textbf{Flowers-102（天花板效应）。} 基线已达到 $99.41\%$ 准确率，几乎没有提升空间。

\textbf{Stanford Dogs（类内变异性）。} Dogs 是唯一 ARR 持续损害性能的数据集。我们假设犬种分类需要与预训练分布显著不同的高度专门化注意力模式。区分密切相关的品种需要关注微妙的、品种特异性的判别区域（耳朵形状、毛色模式、体型比例）。重建损失约束注意力保持接近"良好行为"的分布，可能阻止形成此类精细区分所需的尖锐注意力模式。

\subsection{块放置：深层 vs 中间层}

实验揭示了正则化放置的数据集依赖偏好。一般指南：当任务需要高度局部化的部件级注意力专门化时，重建损失应应用于中间阶段（\texttt{second\_last}）以保留深层灵活性。当任务更全局化时，正则化最后阶段更可取。

\subsection{局限性与未来工作}

\begin{enumerate}
    \item 所有实验使用单一随机种子（42）。需要更大规模的多种子统计评估。
    \item 计算开销（+3.5\,GB 内存，$-18\%$ 吞吐量）在资源受限环境中可能过高。
    \item 仅评估了 Swin Transformer-Base，通用性有待验证。
    \item 一个有��的未来方向是使用预训练模型的注意力作为固定参考，创建更直接的注意力蒸馏形式。
\end{enumerate}

%% ============================================================
\section{结论}
%% ============================================================

我们提出了注意力重建正则化（ARR），一种用于视觉 Transformer 微调的零参数辅助损失，通过惩罚重建注意力分布与观测注意力分布之间的差异来鼓励注意力权重一致性。重建被表述为带闭式单纯形投影的 Tikhonov 正则化逆问题，不添加可学习参数，仅在训练时产生成本。

在七个基准上使用 Swin Transformer-Base 的实验表明，在五个数据集上持续改进，DTD 上增益高达 $+1.55\%$，FGVC Aircraft 上增益高达 $+1.45\%$。全面的消融研究确定 $\alpha = 0.1$ 为稳健的损失权重，\texttt{second\_last} 为普遍有效的放置位置，余弦预热对训练样本少于 $4{,}000$ 的数据集有益。

ARR 与现有正则化技术互补，原则上可与数据增强、标签平滑或知识蒸馏结合。未来工作将聚焦于多种子统计验证、扩展到其他 ViT 架构，以及探索预训练注意力作为固定重建目标。

%% ============================================================
%% 参考文献
%% ============================================================
\begin{thebibliography}{99}

\bibitem{chen2019catastrophic} X.~Chen, S.~Wang, B.~Fu, M.~Long, J.~Wang, Catastrophic forgetting meets negative transfer: Batch spectral shrinkage for safe transfer learning, in: NeurIPS, 2019.

\bibitem{chen2022dearkd} X.~Chen, Q.~Gong, Y.~Zhang, J.~Chen, DearKD: Data-efficient early knowledge distillation for vision transformers, in: CVPR, 2022.

\bibitem{chu2021twins} X.~Chu et al., Twins: Revisiting the design of spatial attention in vision transformers, in: NeurIPS, 2021.

\bibitem{dong2022cswin} X.~Dong et al., CSWin Transformer: A general vision transformer backbone with cross-shaped window self-attention, in: CVPR, 2022.

\bibitem{dosovitskiy2021image} A.~Dosovitskiy et al., An image is worth 16x16 words: Transformers for image recognition at scale, in: ICLR, 2021.

\bibitem{duchi2008efficient} J.~Duchi, S.~Shalev-Shwartz, Y.~Singer, T.~Chandra, Efficient projections onto the $\ell_1$-ball for learning in high dimensions, in: ICML, 2008.

\bibitem{fang2021seed} Z.~Fang et al., SEED: Self-supervised distillation for visual representation, in: ICLR, 2021.

\bibitem{gong2021improve} C.~Gong et al., Improve vision transformers training by suppressing over-smoothing, arXiv:2104.12753, 2021.

\bibitem{hinton2015distilling} G.~Hinton, O.~Vinyals, J.~Dean, Distilling the knowledge in a neural network, arXiv:1503.02531, 2015.

\bibitem{houlsby2019parameter} N.~Houlsby et al., Parameter-efficient transfer learning for NLP, in: ICML, 2019.

\bibitem{hu2022lora} E.~J.~Hu et al., LoRA: Low-rank adaptation of large language models, in: ICLR, 2022.

\bibitem{jia2022visual} M.~Jia et al., Visual prompt tuning, in: ECCV, 2022.

\bibitem{li2018explicit} X.~Li, Y.~Grandvalet, F.~Davoine, Explicit inductive bias for transfer learning with convolutional networks, in: ICML, 2018.

\bibitem{li2019delta} X.~Li et al., DELTA: DEep learning transfer using feature map with attention for convolutional networks, in: ICLR, 2019.

\bibitem{li2023dropkey} B.~Li et al., DropKey for vision transformer, in: CVPR, 2023.

\bibitem{liang2021rdrop} X.~Liang et al., R-Drop: Regularized dropout for neural networks, in: NeurIPS, 2021.

\bibitem{liu2021swin} Z.~Liu et al., Swin Transformer: Hierarchical vision transformer using shifted windows, in: ICCV, 2021.

\bibitem{mobahi2020self} H.~Mobahi, M.~Farajtabar, P.~L.~Bartlett, Self-distillation amplifies regularization in Hilbert space, in: NeurIPS, 2020.

\bibitem{park2022vision} N.~Park, S.~Kim, How do vision transformers work?, in: ICLR, 2022.

\bibitem{raghu2021vision} M.~Raghu et al., Do vision transformers see like convolutional neural networks?, in: NeurIPS, 2021.

\bibitem{szegedy2016rethinking} C.~Szegedy et al., Rethinking the Inception architecture for computer vision, in: CVPR, 2016.

\bibitem{tikhonov1977solutions} A.~N.~Tikhonov, V.~Y.~Arsenin, Solutions of Ill-Posed Problems, Winston, 1977.

\bibitem{touvron2021training} H.~Touvron et al., Training data-efficient image transformers \& distillation through attention, in: ICML, 2021.

\bibitem{vaswani2017attention} A.~Vaswani et al., Attention is all you need, in: NeurIPS, 2017.

\bibitem{wang2021pyramid} W.~Wang et al., Pyramid Vision Transformer: A versatile backbone for dense prediction without convolutions, in: ICCV, 2021.

\bibitem{wang2022minivit} J.~Wang et al., MiniViT: Compressing vision transformers with weight multiplexing, in: CVPR, 2022.

\bibitem{yun2019cutmix} S.~Yun et al., CutMix: Regularization strategy to train strong classifiers with localizable features, in: ICCV, 2019.

\bibitem{zagoruyko2017paying} S.~Zagoruyko, N.~Komodakis, Paying more attention to attention: Improving the performance of convolutional neural networks via attention transfer, in: ICLR, 2017.

\bibitem{zhang2018mixup} H.~Zhang et al., mixup: Beyond empirical risk minimization, in: ICLR, 2018.

\bibitem{zhang2019your} L.~Zhang et al., Be your own teacher: Improve the performance of convolutional neural networks via self distillation, in: ICCV, 2019.

\bibitem{zhou2021deepvit} D.~Zhou et al., DeepViT: Towards deeper vision transformer, arXiv:2103.11886, 2021.

\bibitem{zhuang2020comprehensive} F.~Zhuang et al., A comprehensive survey on transfer learning, Proceedings of the IEEE, 2020.

\end{thebibliography}

\end{document}
